\begin{figure}[H]\centering
% --- (a) przed ruchem: czarny zbije dwa białe w poziomie ---
\begin{subfigure}{0.31\textwidth}\centering
\begin{tikzpicture}[scale=0.66]
\othbg{4}{3}
\othB{0.5}{1.5}\othW{1.5}{1.5}\othW{2.5}{1.5}
\othW{1.5}{0.5}
\othmove{3.5}{1.5}
\draw[-{Latex},thick,black] (3.2,1.5)--(1.05,1.5);
\end{tikzpicture}
\caption{}
\end{subfigure}\hfill
% --- (b) po ruchu: nowy + odwrócone stygną ---
\begin{subfigure}{0.31\textwidth}\centering
\begin{tikzpicture}[scale=0.66]
\othbg{4}{3}
\othB{0.5}{1.5}\othB{1.5}{1.5}\othB{2.5}{1.5}\othB{3.5}{1.5}
\othW{1.5}{0.5}
\othcool{1.5}{1.5}\othcool{2.5}{1.5}\othcool{3.5}{1.5}
\end{tikzpicture}
\caption{}
\end{subfigure}\hfill
% --- (c) whipsaw zablokowany ---
\begin{subfigure}{0.31\textwidth}\centering
\begin{tikzpicture}[scale=0.66]
\othbg{4}{3}
\othB{0.5}{1.5}\othB{1.5}{1.5}\othB{2.5}{1.5}\othB{3.5}{1.5}
\othW{1.5}{0.5}
\othcool{1.5}{1.5}
\othx{1.5}{2.5}
\end{tikzpicture}
\caption{}
\end{subfigure}
\caption{Działanie reguły stygnięcia. (a)~Czarny wykonuje normalne bicie. (b)~Postawiony i~właśnie
odwrócone pionki dostają znacznik stygnięcia (pomarańczowy pierścień) i~są chronione przez jedną
turę. (c)~Biały chciałby natychmiast odbić środkowy pionek z~góry (między swoimi pionkami), ale
ostudzony pionek jest \emph{przezroczysty} -- linia nie domyka żadnego bicia, więc ten ruch jest
nielegalny. Tak reguła eliminuje \emph{whipsaw} (natychmiastowe odwzajemnienie).}
\label{fig:cooldown}
\end{figure}
